% Weekly report to the mentor, 21 November 2026: everything since 15 November (the 14 Nov report).
% Purpose: Lam (author; they/them) keeps track of the project and updates the mentor. Post hoc results
% only, each labelled and placed beside the locked numbers of 14 Nov. No E1-E4 number exists on 21 Nov
% (Phase-1 Colab Session 1 is running); none is written or implied here.
% Build: pdflatex 2026-11-21_weekly_report.tex (twice), or tectonic. No figure and no \input, so the
% .tex builds alone. Standard packages only (same preamble as 2026-11-14_weekly_report.tex), so it also
% builds on Overleaf (upload this file only).
% Source keys used in the % src: comments (line numbers are file lines; CSV line 1 is the header):
%   Files at commit 6e4175d (21 Nov 06:28 +0700, on origin), the head of t3/fm-code-prep when this was written:
%   PR   = protocols.md (6e4175d adds 7 change-log lines after l.503: lines from l.504 on are 7 higher than at fec827e)
%   PH   = results/fm_v1/posthoc/ (post hoc CSVs; every row labelled); G1P = PH G1_comparison.md
%   CX   = docs/research/CONTEXT.md;  DEV = docs/research/FM_DEVIATIONS.md
%   ML   = docs/research/MENTOR_LEARNING_LOG.md (added in 6e4175d)
%   PAP  = docs/report/fm_v1_paper/main.tex; PAPS = docs/report/fm_v1_paper/supplement.tex;
%   PAPR = docs/report/fm_v1_paper/README.md (unchanged since aa722b9)
%   GIT  = git log of branch t3/fm-code-prep (commit times +0700); all commits listed are on origin
% Rounding: half up from the 12-significant-digit form of the stored value (PR l.451-452).
% Not opened for numbers: results/fm_v1/posthoc/pre_rule/ (only its existence is mentioned).
\documentclass[10pt,a4paper]{article}

\usepackage[margin=2cm]{geometry}
\usepackage[T1]{fontenc}
\usepackage{lmodern}
\usepackage{xcolor}
\usepackage{graphicx}
\usepackage{booktabs}
\usepackage{tabularx}
\usepackage{array}
\usepackage{enumitem}
\usepackage{caption}
\usepackage{hyperref}

\definecolor{accent}{HTML}{0550AE}
\hypersetup{colorlinks=true, linkcolor=accent, urlcolor=accent, citecolor=accent}
\setlist{nosep, leftmargin=1.4em}
\captionsetup{font=small, labelfont=bf}
\newcolumntype{Y}{>{\raggedright\arraybackslash}X}
\newcolumntype{R}{>{\raggedleft\arraybackslash}X}
\newcolumntype{P}[1]{>{\raggedright\arraybackslash}p{#1}}
\renewcommand{\arraystretch}{1.15}
\newcommand{\file}[1]{\texttt{\small #1}}
\newcommand{\hash}[1]{\texttt{\footnotesize #1}}
\newcommand{\ask}[2]{\par\smallskip\noindent\textcolor{accent}{\textbf{#1.}} #2\par}

\title{\textbf{Frozen foundation-model encoders for few-shot recognition of rare retinal diseases on OCT}\\[2pt]
  \large Weekly report, 15--21 November 2026: your feedback, the post hoc results and the Phase-1 controls}
\author{Doan Tung Lam}
\date{21 November 2026}

\begin{document}
\maketitle

% ---------------------------------------------------------------------------------------
\section{Summary of the week}

Every locked number of 14 November stands unchanged. You gave feedback twice. The first (17--18 November)
defined ``rare'' by how few labelled images exist and accepted our protocol because it is transparent: no one
can compare with Yoo et al.\ \cite{yoo2021}, whose test split was never published.
The second (20 November) asked for a paper on the reliability and generalisability of the models, not on
MIRAGE beating other backbones, through three controls and with no new model. We computed the agreed post
hoc analyses (P0--P10, R1, R2, G1), each beside the locked numbers. Two findings matter most. A same-picture
audit (P10) found 14 pairs of copies among the 146 rare images, and 4 of the 22 hold-out images have a copy
among the development images; removing them moves the development comparison by less than 0.005. Padding
every image with black to a square (G1) sends almost all external normal scans to rare classes, so external
nine-way results depend on input geometry. We narrowed the novelty claim after finding Sun et al.\ 2024
\cite{sun2024}, and recorded that the public release (146 images) is not the paper's set (119). The rules of
the Phase-1 controls E1--E4 were pushed on 20 November, before any E number; the code and the Colab notebook
followed on 21 November, Step~0 passed and Session~1 is running. No E result exists yet. A first paper draft
was assembled on 20 November. Delivery is 3 December.
% src: locked numbers unchanged: PR l.229-230 (rule 2), PR l.529-532; feedback 1: ML l.10-15, PR l.353-358, l.400-402;
% src: feedback 2: PR l.460-466; post hoc list: PR l.242-259; 14 pairs, 4 of 22, "less than 0.005": PR l.485-496;
% src: PH P10_sensitivity.csv lines 36, 110 (0.115769 -> 0.112920); G1: PR l.411-419; Sun 2024: PR l.31-34, l.429-434;
% src: 146 vs 119: PR l.59-63; E rules pushed 4b4e883 (GIT 20 Nov 22:10), PR l.477-484; code 0578bce, notebook a2c5da2
% src: (GIT 21 Nov 03:13, 03:26); Step 0, Session 1: PR l.499-503; paper: PAPR l.3-6 (assembled 20 Nov; aa722b9, GIT 20 Nov
% src: 21:21); delivery 3 Dec: PR l.189-191

% ---------------------------------------------------------------------------------------
\section{What the mentor said and what changed}

\textbf{Feedback 1 (17--18 November).} ``Rare'' should be judged by how few labelled images exist, not by
how common the disease is. The methods are simple. Later (18 November): the protocol and the existing results
are acceptable because they are transparent; on comparability with Yoo et al., ``even the reviewer cannot
identify clearly what it is because no one can compare'' (nobody has their test split).
% src: ML l.12-15; PR l.353-354 (mentor answered 17 Nov), l.400-402 (18 Nov, quote verbatim)

\smallskip
\noindent{\small
\begin{tabularx}{\linewidth}{@{}P{3.1cm} Y P{2.6cm}@{}}
\toprule
Request & What we did & Where recorded \\
\midrule
Define ``rare'' by label scarcity &
Definition B1 (Lam and you, 18 Nov): \emph{rare} = at most 5\% of the dataset's labelled images (at most 3\%
also stated); \emph{few-shot} = example pool of at most 30. Tiers: A rare and few-shot; B rare, pool over 30
(capped); C few-shot, over 5\%; D common. Yoo's five classes are tier A (pools of 19--27; class totals
0.55--0.77\% of the 4,146 images used); OCTID CSC and MH tier D (102 each, 17.8\% of 572); OCTDL MH tier B
(1.9\% of images; 28 of 821 patients, 3.4\%). External sets are worded as agreement with published labels. &
\file{protocols.md} \S8, 18 Nov; \hash{39c3d73}, \hash{dcbb323} \\
The methods are simple &
No method added to this study: the design stays as locked, and the protocol forbids changes after results.
Fine-tuning (LoRA) moved to the next phase, on data where image size does not predict the class; Phase 2
(after 3 Dec) studies few-shot adaptation methods. &
\file{protocols.md} \S6, \S8; \hash{dcbb323}, \hash{0b051c4} \\
Transparent is enough; ``no one can compare'' &
No forced comparison: both locked sets stay, no new split, no new number. ``A hold-out matching their class
counts'', never ``their test set''; Inception-v3 is ``the backbone family of Yoo et al.''; their 92.1 and 93.9
are context only. Provenance recorded; novelty narrowed (Section~\ref{sec:prov}). &
\file{protocols.md} \S1, \S3, \S8; \hash{ca0d26d}; \file{CONTEXT.md} 19 Nov \\
Our follow-up (problem B5): hold-out leakage &
Same-picture audit P10: rule 18 Nov, result 21 Nov (Section~\ref{sec:p10}). &
\hash{dcbb323}, \hash{01c3b7c}, \hash{917ab3b} \\
\bottomrule
\end{tabularx}}
% src: B1 and tiers: PR l.359-368 (commit dcbb323, GIT); 17.8% of 572 and 1.9% of 2,064: PR l.354-356 (commit 39c3d73);
% src: OCTDL MH 1.9% by images, 28 of 821 patients = 3.4%: PR l.380-381; class totals 23-32 = 0.55-0.77% of 4,146: PR l.364-366;
% src: wording: PR l.380-384; methods simple: ML l.13; no outcome-driven change: PR l.155-156; fine-tuning next phase:
% src: PR l.392-393 (dcbb323); Phase 2: PR l.474-476 (0b051c4); comparability: PR l.425-431 (ca0d26d), PR l.56-58;
% src: CX l.127 (CONTEXT.md decisions row 2026-11-19); P10 for B5: PR l.392 (dcbb323); GIT 01c3b7c (19 Nov), 917ab3b (21 Nov)

\medskip
\noindent
\textbf{Feedback 2 (20 November).} The line from GAN augmentation to foundation models, and the findings on
source confounding and external degradation, have real scientific value. A paper should focus on the
reliability and generalisability of the models, not on MIRAGE beating other backbones. Three things to do:
(1) control for image size and framing and check whether MIRAGE keeps its advantage (the most critical);
(2) error analysis on OCTID and OCTDL, especially normal scans called other diseases; (3) sensitivity of
MIRAGE and RETFound to preprocessing. No more models and no more GAN work now.
% src: PR l.460-466; ML l.34-38

\smallskip
\noindent{\small
\begin{tabularx}{\linewidth}{@{}P{3.1cm} Y P{2.6cm}@{}}
\toprule
Request & What we did & Where recorded \\
\midrule
Reliability and generalisability, not ranking &
Two phases. Phase 1 (to 3 Dec): image size, framing and source, and their effect on the models; one paper on
3 Dec holds the locked results and these controls (delivery moved back from 30 Nov). The draft introduction
carries a reframing sentence for Lam to approve; the full reframing is planned for 28--30 Nov. &
\file{protocols.md} \S8; \hash{0b051c4}; paper \hash{aa722b9} \\
(1) Size and framing & E1, all four models: one main control, one success rule (Section~\ref{sec:phase1}). &
rules \hash{4b4e883}, code \hash{0578bce}, notebook \hash{a2c5da2} \\
(2) External errors & E2 (laptop): where external normal scans go, which scans the models fail together. &
\hash{4b4e883} \\
(3) Preprocessing & E3: a fixed grid of one-step changes to the MIRAGE and RETFound recipes. & \hash{4b4e883} \\
No new model, no GAN & None added; the rules allow no new model, no GAN image and no new download. &
\hash{4b4e883} \\
(Lam's addition) & E4: how easily the four OCT sources can be told apart. & \hash{4b4e883} \\
\bottomrule
\end{tabularx}}
% src: two phases: PR l.467-476 (0b051c4, GIT 20 Nov 21:00); reframing TODO: PAP l.100; 28-30 Nov: PR l.284;
% src: 30 Nov -> 3 Dec: PR l.189-191; E1-E4 and their mapping to M1-M3: PR l.514-519, l.594-597 (4b4e883, GIT 20 Nov 22:10);
% src: no new model/GAN/download: PR l.518-519; code and notebook: GIT 0578bce, a2c5da2; PR l.639

% ---------------------------------------------------------------------------------------
\section{Results added since 15 November (all post hoc)}
\label{sec:results}

Every result here was computed after the locked results of 14 November, sits beside them under a dated post
hoc label and never replaces them. Each rule was committed and pushed before its numbers, except P0, P6 and
P7, first looks that are labelled so. Everything is released, favourable or not.
% src: PR l.227-237 (rules 1-4), l.333-342 (P6 label; rule 1 met 16 Nov 17:56, commit 1eb3964), l.229 (P0, P7 label)

\subsection{Development spread over support draws (R2)}
R2 adds the spread that protocol section~5 promised and the locked notebook did not write: rule committed
19 November (\hash{65fe98c}), computed 20 November after the results were known. Gate passed: every mean
equals the locked nomination value within $10^{-9}$. The paper uses it as Table~I.
% src: PR l.257 (R2 rule, gate 1e-9, label, Table 1), l.458-459 (computed 20 Nov after 65fe98c and 0ceeb56; gate passed);
% src: PH R2_development_spread.csv column posthoc ("rule timestamped by commit 65fe98c (19 Nov 2026), computed 2026-11-20")

\medskip
\noindent{\small
\begin{tabularx}{\linewidth}{@{}P{3.6cm} R R R R R R >{\raggedleft\arraybackslash}p{2.3cm}@{}}
\toprule
Development macro recall & $K=1$ & $K=2$ & $K=4$ & $K=8$ & Mean $K$1--8 & SD & 5th--95th pct \\
\midrule
\textbf{MIRAGE} (nominee) & 0.495 & 0.615 & 0.712 & 0.777 & \textbf{0.650} & 0.028 & 0.603 to 0.688 \\
% src: PH R2_development_spread.csv lines 12-16
DINOv2 & 0.407 & 0.497 & 0.616 & 0.711 & 0.558 & 0.028 & 0.511 to 0.606 \\
% src: PH R2_development_spread.csv lines 7-11
Inception-v3 (reference) & 0.409 & 0.499 & 0.581 & 0.647 & 0.534 & 0.029 & 0.494 to 0.583 \\
% src: PH R2_development_spread.csv lines 17-21
RETFound & 0.318 & 0.375 & 0.442 & 0.489 & 0.406 & 0.024 & 0.367 to 0.441 \\
% src: PH R2_development_spread.csv lines 2-6
\midrule
MIRAGE $-$ Inception-v3 & $+$0.086 & $+$0.116 & $+$0.132 & $+$0.130 & $+$\textbf{0.116} & 0.037 & $+$0.069 to $+$0.185 \\
% src: PH R2_development_spread.csv lines 32-36
DINOv2 $-$ Inception-v3 & $-$0.002 & $-$0.002 & $+$0.036 & $+$0.064 & $+$0.024 & 0.034 & $-$0.025 to $+$0.086 \\
% src: PH R2_development_spread.csv lines 27-31
RETFound $-$ Inception-v3 & $-$0.090 & $-$0.124 & $-$0.139 & $-$0.159 & $-$0.128 & 0.035 & $-$0.186 to $-$0.070 \\
% src: PH R2_development_spread.csv lines 22-26
\midrule
\textit{Size-only reference (locked)} & \textit{0.578} & \textit{0.550} & \textit{0.560} & \textit{0.589} & \textit{0.569} & & \\
% src: G1P l.28-32, column "Size only" (locked size-only reference)
\textit{Chance} & \textit{0.200} & \textit{0.200} & \textit{0.200} & \textit{0.200} & \textit{0.200} & & \\
% src: PR l.89-90 ("Guessing gives 0.20")
\bottomrule
\end{tabularx}}

\smallskip
\noindent{\small 124 development queries, five folds pooled; mean over 100 support draws, every model on the
same supports. Paired rows: difference from Inception-v3 on the same draw. SD and percentiles (for the mean
over $K$) describe variation over draws on fixed images; they are not confidence intervals.}
% src: PR l.257 (unit: 124 queries, folds pooled), l.145-147 (spread is not a confidence interval)

\medskip
\noindent
The 5th percentile of MIRAGE minus Inception-v3 is above zero at every $K$ ($+0.001$, $+0.015$, $+0.052$,
$+0.071$). DINOv2 minus Inception-v3 includes zero at every $K$. RETFound minus Inception-v3 has its 95th
percentile below zero at every $K$ ($-0.0003$ at $K=1$). The size-only reference beats every model at $K=1$.
% src: PH R2_development_spread.csv p05 lines 32-35 (0.000848, 0.015124, 0.051972, 0.071306); DINOv2 p05/p95 lines 27-31
% src: (all p05 < 0 < p95); RETFound p95 lines 22-26 (max -0.000319 at K=1); size-only 0.578 at K=1 (G1P l.28) > 0.495 (R2 line 12)

\subsection{Post hoc analyses P0--P9 and R1}

\noindent{\small
\begin{tabularx}{\linewidth}{@{}l P{3.0cm} Y P{1.9cm}@{}}
\toprule
& Analysis & Key result & Placement \\
\midrule
P0 & Recall per disease (first look) & MIRAGE is not the best model for macular hole: 0.540, against 0.612 for
DINOv2 and 0.623 for Inception-v3 (development, mean $K$1--8). & Supplement \\
% src: PH P0_per_class_recall.csv lines 15, 9, 21 (mean_k1_2_4_8); label "rule fixed after first look"; placement PR l.447-448
P1 & Exact bounds; leave one image out & MIRAGE called 0 of 169 patient-disjoint Kermany normals rare: exact
95\% upper bound 0.022. With any one image left out, MIRAGE minus Inception-v3 (mean $K$) stays within
0.109--0.122 on development (all: 0.116) and 0.173--0.232 on the hold-out (all: 0.197). & Supplement \\
% src: PH P1_nine_way_exact_bounds.csv line 116 (cp95_high 0.0216); PH P1_leave_one_out_summary.csv lines 2-3
P3 & Clinic-to-clinic macular hole (outside the locked study) & OCTID examples, OCTDL queries, two-way (chance
0.50): DINOv2 0.858, RETFound 0.675, MIRAGE 0.550 (its recall on OCTDL normals 0.130), Inception-v3 0.500.
Reverse: MIRAGE 0.691. Worded as agreement between two macular-hole labels. & Supplement \\
% src: PH P3_clinic_to_clinic_mh.csv lines 21, 6, 36, 51 (image units, mean_k1_2_4_8), line 36 recall_NORMAL_mean 0.130,
% src: line 46 (reverse, MIRAGE 0.691); wording PR l.247; placement PR l.384-385
P4 & Size-right against size-wrong queries & The size-only rule names 0.594 of development query cells
correctly; every model scores higher on those (MIRAGE 0.663 against 0.646; Inception-v3 0.543 against
0.509). A description, not proof. & Main text \\
% src: PH P4_size_right_vs_wrong.csv line 42 (0.594, 29441/49600), lines 26, 31, 36, 41; PR l.248; placement PR l.446-447
P5 & Where normal scans go; ``not sure'' cut-off & Of 332 OCTDL normals, MIRAGE calls 26 normal, 276 a common
disease and 30 rare; 295 fall below its ``not sure'' cut-off, set on Kermany normals only. & Supplement \\
% src: PH P5_not_sure_cutoff.csv line 10 (normals 332, not_sure 295, false_rare 30, called_normal 26, called_common 276); PR l.249
P8 & CSC against MH on development & MIRAGE minus Inception-v3 $-0.083$ (5th--95th $-0.164$ to $+0.010$),
against $+0.026$ on the OCTID two-way test; size-only 0.968, above every model. & Main text \\
% src: PH P8_csc_vs_mh_development.csv line 31 (-0.0834, p05 -0.1636, p95 0.0104, locked OCTID 0.0261), line 26 (0.968)
P9 & Rare against normal, AUROC, no cut-off & MIRAGE: OCTID CSC+MH 0.781 (95\% 0.733--0.824); OCTDL MH 0.357
(0.261--0.468), so this score ranks OCTDL scans the wrong way; DINOv2 on OCTDL 0.978. & Supplement \\
% src: PH P9_rare_vs_normal_auroc.csv lines 12, 16, 15; reading of MIRAGE's OCTDL value: PR l.390-392
R1 & External disease images beside each false-rare rate & The MIRAGE head calls 41 of 102 OCTID CSC, 11 of
102 OCTID MH and 3 of 39 OCTDL MH images any rare class: a low false-rare rate can mean it rarely says
``rare''. & Beside every false-rare rate \\
% src: PH R1_rare_beside_false_rare.csv lines 11-13; PR l.259
\bottomrule
\end{tabularx}}

\smallskip
\noindent
Also in the supplement: P2 (October's trained heads on OCTID and OCTDL; gates passed, 17, 18 and 19 of 22
reproduced), P6 (prototype against logistic probe) and P7 (MIRAGE+DINOv2 ensemble: development 0.695 against
0.650, hold-out 0.718 against 0.722; one line, never called a result).
% src: P2 gates PR l.346; P6 PR l.250; P7 PR l.251; placement PR l.446-448

\subsection{G1: pad-to-square input check}
G1 re-extracted MIRAGE and Inception-v3 with every image padded with black to a square, and scored the locked
episodes and roster. It ran on Colab on 18 November; unpadded controls and the 1,032 square images matched the
locked features within $1.43\times10^{-6}$. G1 goes to the supplement; the main text keeps the locked 26 of
332 OCTDL normals called normal, with no claim about its cause. Lam's decisions of 19 November fix the wording:
% src: PR l.255 (rule), l.403-410 (ran 18 Nov 07:02 UTC; 1,032 square images within 1.43e-6), l.411-414 (placement, 26 of 332)

\smallskip
\noindent{\small
\begin{tabularx}{\linewidth}{@{}Y P{6.2cm}@{}}
\toprule
Allowed & Not allowed \\
\midrule
Development and hold-out recall moved by at most 0.023 at $K=1$--8: a near-null test, because the rare inputs
are already near square. &
That squashing did or did not cause the OCTDL result; robustness to aspect ratio. \\
With black padding, external normal scans go almost entirely to rare classes. Called NORMAL: MIRAGE OCTID
202 $\rightarrow$ 0, OCTDL 26 $\rightarrow$ 1; Inception-v3 182 $\rightarrow$ 8, 66 $\rightarrow$ 13. &
Padded external ``gains'' as recognition; ``rare is a not-Kermany bin''; ``bars caused it''. \\
The rare images are letterboxed and padding gives external scans that layout: external nine-way results
depend on input geometry. &
Any reading of $K=16$; any MIRAGE-versus-Inception-v3 ranking from padded scores. \\
P9 caveat: the locked OCTID AUROCs hold ``under the published input path; with black bars added, the same
score ranks OCTID the other way'' (MIRAGE 0.781 $\rightarrow$ 0.198). &
Anything about what the features contain. \\
\bottomrule
\end{tabularx}}
% src: claims from PR l.414-424; largest change 0.023: G1P l.28-31, l.40-43 (MIRAGE hold-out K=2 -0.023);
% src: called NORMAL: G1P l.67-68, l.73-74; P9 padded: G1P l.84 (OCTID CSC+MH 0.781 -> 0.198); optional G1 sentence: PR l.423-424

\subsection{P10: same-picture audit of the 146 rare images}
\label{sec:p10}
Yoo et al.\ took the rare images from web searches, so one published picture could sit on both sides of a
split. Their source list gives links per disease, not per file, so no file could be mapped to a link. A
keypoint screen (SIFT with RANSAC; score = inlier count) scored all 10,585 pairs. The cut-off, 18 inliers, is
the largest score among 2,000 rare--Kermany pairs, which cannot show the same picture; re-cropped synthetic
copies reached it for 146 of 146 images. Lam marked 55 items blind to label, split and score (the 35 pairs at
or above the cut-off, 10 synthetic copies, 10 decoys), and the marks were committed (\hash{abb1770}) before the
key was joined. All copies were marked ``same picture'' and all decoys ``different''. Of the 35 pairs: 14
``same picture'', 1 ``same case, other slice'' (MacTel 09/23), 1 ``unsure'' (CSC 26/27; Lam's note: probably a
follow-up visit) and 19 ``different''. The locked duplicate screen had flagged none of the 14. \textbf{Four of the 22
hold-out images have a same-picture copy in development:} MacTel/19, RP/17, RP/20 and Stargardt/03.
% src: PR l.254 (rule), l.485-491; PH P10_screen_summary.csv lines 2, 3, 5, 6, 8, 9, 12, 13; P10_source_links.csv ("no file
% src: named: S1 lists links per disease"; PAP l.233-234); P10_review_key.csv kinds (35 pair, 10 synthetic, 10 decoy);
% src: P10_review_marks.csv joined with the key (synthetic 10 same picture; decoy 10 different; pairs 14/1/1/19); GIT abb1770;
% src: PH P10_matched_pairs.csv lines 6, 10, 12, 13-14 (hold-out images)

\medskip
\noindent{\small
\begin{tabularx}{\linewidth}{@{}P{4.3cm} P{3.1cm} R R R@{}}
\toprule
P10 sensitivity (mean over $K$1--8) & Excluded & Development: MIRAGE $-$ Inception-v3 (5th pct) & Hold-out:
MIRAGE & Hold-out: MIRAGE $-$ Inception-v3 (5th--95th) \\
\midrule
Locked (no exclusion) & none & $+$0.116 ($+$0.069) & 0.722 & $+$0.197 ($+$0.042 to $+$0.343) \\
% src: PH P10_sensitivity.csv lines 36, 54, 75
Without ``same picture'' or ``same case'' & 1,196 query-episodes; 4 hold-out images & $+$0.113 ($+$0.064)
& 0.751 & $+$0.240 ($+$0.089 to $+$0.407) \\
% src: PH P10_sensitivity.csv lines 110, 128, 149 (excluded column: 1196 query-episode exclusions per arm; 4 hold-out images)
Also counting ``unsure'' & 1,378 query-episodes; 4 hold-out images & $+$0.113 ($+$0.064) & 0.751 & $+$0.240
($+$0.089 to $+$0.407) \\
% src: PH P10_sensitivity.csv lines 184, 202, 223
\bottomrule
\end{tabularx}}

\smallskip
\noindent
Agreed wording: ``a keypoint screen and a blind review found 14 same-picture pairs; removing them changes the
development comparison by less than 0.005''; never ``no shared patients''. E1--E4 use the locked episodes, so
these pairs sit inside them; E1 will be read beside this sensitivity, not re-run without them.
% src: PR l.254 (exclusion rule), l.494-498 (wording, E1 read beside); 0.116 - 0.113 = 0.0028 (PH P10_sensitivity.csv lines 36, 110)

\subsection{Dataset provenance, layout and novelty}
\label{sec:prov}
\textbf{Provenance} (file identity only). Mendeley btv6yrdbmv version~1 (27 March 2020) holds 119 images, the
total in Yoo et al. Version~2 (29 October 2020), which we use, holds 146: the 118 version-1 files unchanged plus
28 new files (27 new names and a different picture under MacTel/20.PNG). 21 of the 22 hold-out images are in
version~1. RP/20.PNG is a version-2 addition and, P10 found, a copy of the development image RP/03. Yoo et al.\
give Stargardt 15+4 and RP 12+4, but the version-1 folders hold RP 19 and Stargardt 16 (as does Sun et al.), so
the two names may be swapped in one of them; per-class results for these two carry a footnote, and macro
recall is unaffected. The mapping file sits in \file{results/fm\_v1/posthoc/pre\_rule/}, released outside any
rule.
% src: PR l.59-67 (counts, dates, footnote); PR l.490-491 (RP/20 a version-2 addition, a copy of RP/03); PR l.437-440 (pre_rule);
% src: "the two names may be swapped in one of them": PAP l.509

\smallskip
\noindent
\textbf{Layout.} In the locked geometry summary, the median share of flat border per rare class is 0.34--0.50
(development) and 0.22--0.56 (hold-out), against 0 (at most 0.002) for Kermany, OCTID and OCTDL. E1's main
control crops these borders away. One deviation from a locked text is recorded (D1): A2 and A3 call the OCTID
two-way test the one ``where image size is constant''; it is not, and since every OCTID image is larger than
every benchmark CSC and MH image, the size-only rule sends all its queries to CSC.
% src: PR l.68-71; E1 content box: PR l.565-569, l.594; DEV l.10 (D1, 20 Nov; "sends all 204 queries to CSC and scores 0.50")

\smallskip
\noindent
\textbf{Novelty.} Sun et al.\ \cite{sun2024} already drew five-way $K$-shot tasks ($K = 1$, 3, 5) from the same
rare images, with a trained network; it is not a numeric comparator (other pool and training). Our claim is
narrowed to ``a pre-registered, shortcut-audited evaluation of frozen foundation models on a public rare-OCT
benchmark'': frozen encoders at 1 to 16 images per disease on identical supports, how often normal scans from
other sources are called rare, and what image size and source alone explain. Lam's positioning, to be
confirmed with you at the 23 November review.
% src: PR l.31-37 (narrowed 19 Nov), l.432-434, l.674

% ---------------------------------------------------------------------------------------
\section{Phase 1 (to 3 December): the controls E1--E4}
\label{sec:phase1}

One agent drafted the rules, a second rebutted them (19 findings, 6 must-fix), and Lam fixed the open
choices. No E number exists yet.
% src: PR l.477-484 (Opus 5.5 draft, Fable 5.2 rebuttal: 19 findings, 6 must-fix; Lam's choices)

\medskip
\noindent{\small
\begin{tabularx}{\linewidth}{@{}P{1.75cm} Y Y P{2.35cm}@{}}
\toprule
Request & Design & Success rule or criteria (fixed 20 Nov) & Status (21 Nov) \\
\midrule
\textbf{E1} (M1): size and framing &
Main control E1-M, one render for all four models: crop the flat borders, take the centred square, resize to
$256\times256$; each model then applies its locked transform. Descriptive variants E1-L (resolution floor)
and E1-C (crop only); references without a model (size-only, content size, $16\times16$ thumbnail). &
Per draw, mean over $K=1,2,4,8$ of MIRAGE E1-M $-$ Inception-v3 E1-M development macro recall. 5th percentile
over the 100 paired draws above 0: MIRAGE keeps its advantage under this control; otherwise not; a failed
control: ``not evaluable''. Verdict from E1-M alone, in sentences fixed in advance; hold-out for direction
only. &
Session~1 (E1-M; E1-L for MIRAGE and Inception-v3): running. Session~2: the rest of E1. \\
\textbf{E2} (M2): external errors &
Laptop, on locked and E1 nine-way outputs: where the 206 OCTID and 332 OCTDL normals go; how many models fail
the same scans; locked $\rightarrow$ E1-M moves; seven image properties of failed scans; OCTDL per patient; a
blinded contact sheet marked by Lam (non-clinician). &
Descriptive only: no test, no $p$-value, no cause, no judgement of a dataset's labels; any association quoted
for all four models. &
Laptop, after Session~2. \\
\textbf{E3} (M3): preprocessing &
One step changed at a time. MIRAGE: centre crop, bicubic, no min--max stretch, 256-pixel input (G1's padding
beside, known in advance). RETFound: squash, black pad, bilinear, min--max stretch, mean and SD 0.5, 128-pixel
short edge. &
MIRAGE: every variant $-$ Inception-v3 has its 5th percentile (mean over $K$) above 0: ``kept within this
grid''. RETFound: every variant below Inception-v3 (95th percentile) and below MIRAGE (5th percentile of MIRAGE
$-$ variant above 0). False-rare counts inside the locked 95\% intervals: ``stable''. Input preprocessing
only. &
Session~1 (all of E3): running. \\
\textbf{E4} (Lam): source separability &
Kermany (1,000 test), Yoo's 146, OCTID 410, OCTDL 371. Main: a few-shot source test on the three normal-only
pairs, locked and E1-M features. Secondary: AUROC per image property; grouped logistic probe. &
Fixed readings: AUROC $\ge 0.95$ ``near-perfectly separable'', 0.80--0.95 ``strongly'', 0.60--0.80 ``partly'',
below 0.60 ``hardly''. No claim that models use the source in the rare task. &
Laptop, after Session~2. \\
\bottomrule
\end{tabularx}}
% src: E1: PR l.594 (E1-M, E1-L, E1-C, references, success rule, fixed sentences, hold-out); E2: PR l.595 (seven common
% src: properties; Lam marks the sheet: PR l.480-481); E3: PR l.596 (variants, criteria 1-3, scope); E4: PR l.597 (sources and
% src: sizes, main measure, readings, not allowed); M1-M3: PR l.514-519; sessions: PR l.599-603; status: PR l.499-503, l.639

\medskip
\noindent
\textbf{Integrity.}
\begin{itemize}
  \item \emph{Rules before numbers.} The rules were pushed (\hash{4b4e883}, 20 November 22:10) before any E
        number; the scoring code checks that commit and hashes the rule text. No analysis rule is added after
        this push; a later edit may only fix a bug that stops a run.
  \item \emph{Written knowing the results.} Two cells were known in advance and are labelled so (E3's padded
        MIRAGE is G1; the locked values beside every E number are R2's). Hence one main control and one
        success rule; E1's variants, E3 and E4 cannot overturn or rescue the E1 verdict.
  \item \emph{Colab checks.} For each model, unchanged originals through the locked path must match the locked
        features within $10^{-4}$ (all 5,927 images; the first 20 for Inception-v3), or its numbers are ``not
        interpreted''. A failed step may be rerun once, unchanged.
  \item \emph{Scored once.} Before any E1 result existed, Lam committed to run Session~2 right after
        Session~1, whatever Session~1 shows. E1--E4 are scored once, after Session~2; Session~1 outputs get
        integrity checks only.
  \item \emph{Images stay in Drive,} with a SHA-256 receipt per archive; only CSVs and generated pages enter the
        repository.
\end{itemize}
% src: PR l.523-528 (rule 1), GIT 4b4e883 20 Nov 22:10 +0700; PR l.533-537 (rule 3); PR l.544-545 (rule 5);
% src: PR l.585-590 (Colab checks, 5,927 images, 1e-4), l.546-549 (rule 6); PR l.499-503 (scored once); PR l.550-552 (rule 7)

\noindent
\textbf{Status.} Rules \hash{4b4e883} (20 Nov 22:10); code \hash{0578bce} (21 Nov 03:13); notebook
\file{colab\_fm\_control.ipynb} pinned to it (\hash{a2c5da2}, 03:26), whose laptop dry run reproduces R2.
On 21 November Step~0 passed (MIRAGE, RETFound and DINOv2 reachable, after a token-access fix) and
Session~1 started (planned about 3~h); Session~2 (about 1.5~h) follows, then E1--E4 are scored once.
% src: GIT (commit times +0700); PR l.639 (notebook pinned, dry run reproduces R2); PR l.499-503 (Step 0, HF_TOKEN fix,
% src: Session 1, Session 2 commitment; recorded in fec827e); PR l.601-603 (Session 1 about 3 h, Session 2 about 1.5 h)

% ---------------------------------------------------------------------------------------
\section{Paper status}

The draft (\file{docs/report/fm\_v1\_paper/}, IEEE conference template) was assembled on 20 November from the
reviewed section drafts. Its tables are generated from the locked and post hoc files, and every numeric
sentence has a source comment. The P10 text and the Phase-1 methods were added on 21 November
(\hash{6e4175d}), after two rounds of agent review. \emph{Done}: text and numbers complete, awaiting Lam's
final read; \emph{draft}: text with open items; \emph{TODO}: not written.
% src: PAPR l.1-17; PR l.626 ("rebuilt, never hand-edited numbers"); GIT 6e4175d (21 Nov 06:28, "Paper: add P10 and the E1-E4
% src: Methods"; message: "Reviewed with agent-loop (2 rounds ...)"); PAP l.15-24 (sources for the text added on 21 Nov)

\medskip
\noindent{\small
\begin{tabularx}{\linewidth}{@{}P{6.3cm} P{1.5cm} Y@{}}
\toprule
Part & Status & Open \\
\midrule
Title, abstract, I Introduction & Draft & Lam approves; reframing sentence; a shortcut-learning citation (none
among the tracked sources) \\
II Related work & Draft & Title and DOI of Sun et al.\ 2024 unverified \\
III Methods, including post hoc and Phase-1 methods & Done & -- \\
IV Results: locked (R2 as Table~I) and post hoc (P4, P8, P10) & Done & Placement (Lam) \\
IV Results: Phase-1 controls, inside the one dated post hoc subsection & \textbf{TODO} & After E1--E4 are
scored; the default main-text tables (E1 success, E1 nine-way, E2 summary, E2d) appear when their files exist \\
V Discussion; VI Limitations (includes P10) & Draft & Phase-1 parts TODO \\
VII Next phase & Draft & Whether to name the planted-cue test (Lam) \\
VIII Conclusion & \textbf{TODO} & Draft only, for Lam \\
Declarations & Draft & Zenodo time; data and code archive; AI use; ethics, funding; venue not chosen \\
Supplement & Draft & Locked S1--S7; P0--P3, P5--P7, P9, G1, P10; E1--E4 placeholders \\
\bottomrule
\end{tabularx}}
% src: PAP section lines: title/abstract TODO l.57, l.64; I l.72 (TODOs l.85, l.100); II l.122; III l.139-251 (post hoc l.226-234,
% src: Phase-1 l.236-251; D1 l.211-212); IV l.253-434 (post hoc subsection l.352, Phase-1 subsubsection l.387, TODO l.389;
% src: E tables guarded by \havetable: e1d_success, e1g_nine_way, e2_destinations_summary, e2d_octdl_locked l.391-434);
% src: V l.436-490 (TODO l.489); VI l.492-519; VII l.521-522 (planted-cue TODO); VIII l.525-526 (TODO); declarations l.530-550;
% src: PAPR l.47-53 (open items, incl. Sun 2024 bib fields); PAPS l.18-242 (S-A S1-S7, S-B P0-P10, S-D E1-E4)

% ---------------------------------------------------------------------------------------
\section{Calendar to 3 December, and Phase 2}

Dates are the latest dates; when a step finishes early, the next one starts at once. Lam does the slow steps,
the agents verify and draft at the same time, and your feedback is welcome whenever it comes.
% src: PR l.193-196 (pull-forward rule; parallel work; mentor not a gate)

\medskip
\noindent{\small
\begin{tabularx}{\linewidth}{@{}l Y P{3.0cm}@{}}
\toprule
Dates & Work & Status (21 Nov) \\
\midrule
16--18 Nov & P0--P9, R1; G1 run on Colab and scored & Done \\
19--21 Nov & P10; R2; paper draft; Phase-1 rules, code and notebook & Done \\
21--23 Nov & E1 and E3 on Colab (Sessions 1 and 2); 23 Nov: review with you (positioning) & Session~1 running \\
24--26 Nov & E1--E4 scored once; E2 and E4 on the laptop; one verdict per experiment; figures and tables & -- \\
\textbf{27 Nov} & \textbf{All numbers, tables and figures frozen} & -- \\
28--30 Nov & Paper reframed on reliability and generalisability; discussion, limitations, next phase & -- \\
1--3 Dec & Evidence audit (1 Dec), PDF freeze (2 Dec), \textbf{delivery (3 Dec)} & -- \\
\bottomrule
\end{tabularx}}
% src: PR l.272-285 (calendar rows); PR l.186 (P0-P9, G1 done 18 Nov; R2 20 Nov); PR l.403-409 (G1 ran 18 Nov);
% src: PR l.485-503 (P10 done; Session 1 started 21 Nov, ahead of the 22-23 Nov slot); PR l.37, l.433-434 (23 Nov review)

\smallskip
\noindent
No new analysis rule after the Phase-1 push; feedback after the freeze changes wording only, never a number.
\textbf{Phase 2 (after 3 December)} returns to the original direction: few-shot adaptation methods across
foundation models, datasets and rare-label tiers, with the Phase-1 shortcut audit as a standard check. It may
start once the numbers are frozen; its results stay out of this paper except one labelled paragraph.
Fine-tuning runs in order, under its own protocol: a test where image size does not predict the disease
(patient IDs if possible); light adaptation (LoRA or the last blocks) of MIRAGE and RETFound in the same
$K$-shot episodes; full fine-tuning only if that gains clearly on development and holds on OCTID and OCTDL.
Decided for the next version: remove the letterbox and rescore. Open (Lam): a planted-cue test; DINOv3.
% src: PR l.196-199 (gates); PR l.474-476 (Phase 2); PR l.198-200 (starts when frozen; results out of this paper);
% src: PR l.313-320 (fine-tuning order); PR l.308-309, l.435-436 (letterbox removal, decided 19 Nov); PR l.310-311
% src: (planted-cue, Lam's choice pending); PR l.393-394 (DINOv3: v2 or next phase decided after 3 Dec)

% ---------------------------------------------------------------------------------------
\section{Open questions for you}

\ask{Q1}{Positioning. Is ``a pre-registered, shortcut-audited evaluation of frozen foundation models on a
public rare-OCT benchmark'', written around reliability and generalisability rather than which model wins, the
framing you want for the 3 December paper? We would like to confirm it at the 23 November review.}
% src: PR l.33-37, l.432-434; PR l.460-463; PAP l.98-100

\ask{Q2}{The hold-out after P10. Four of the 22 hold-out images have a copy in development. The locked
22-image numbers stay; the draft gives the values without the four (MIRAGE minus Inception-v3 $+0.240$,
against the locked $+0.197$) in the main-text post hoc paragraph and names the copies in the Limitations. Is
that enough, or should the hold-out table itself carry a labelled row without the four images?}
% src: PR l.489-494, l.229-230 (rule 2); PH P10_sensitivity.csv lines 75, 149; PAP l.363-366 (post hoc paragraph), l.500 (Limitations)

\ask{Q3}{Placement. By default the E1 success table, the E1 nine-way audit and an E2 summary go in the main
text's post hoc subsection, with the E2 table of OCTDL image properties that the E2 rule fixes for the main
text; everything else goes to the supplement. One optional main-text sentence would point to G1. Do you agree
with this placement and with that sentence?}
% src: PR l.610-612 (default placement), l.595 (E2 rule: "The main-text table is fixed now: OCTDL normals, locked features, the
% src: seven common properties x four arms"), l.608; PAP l.425-432 (e2d_octdl_locked); PR l.423-424 (optional G1 sentence)

\ask{Q4}{A second rater. Lam alone, a non-clinician, marked the P10 pairs and will mark the E2 contact sheet
(24--26 November). Do you know a clinician or a second rater who could mark the same sheet blind? Otherwise it
stays a stated limitation.}
% src: PR l.480-481 (contact sheet marked by Lam alone, as a non-clinician), l.595 ("One rater, not a clinician (Part E may
% src: add a second)"); PAP l.494 (one non-clinician rater, P10); calendar PR l.282

\medskip
\noindent\rule{\linewidth}{0.4pt}\\
{\footnotesize \textbf{Record.} Every rule, decision and date above is in
\href{https://github.com/Malnoad/Few-shot-rare-diseases/blob/t3/fm-code-prep/protocols.md}{\file{protocols.md}},
section~8. Commits since 15 November, all pushed to \file{t3/fm-code-prep}: post hoc plan \hash{1eb3964}
(16 Nov); P9 and G1 code \hash{39c3d73}, P10 rule and B1 \hash{dcbb323} (18 Nov); G1 outputs \hash{d25e152},
19 Nov decisions \hash{ca0d26d}, R2 rule \hash{65fe98c} (19 Nov); R2 \hash{24da8f3}, two-phase plan
\hash{0b051c4}, paper \hash{aa722b9}, Phase-1 rules \hash{4b4e883} (20 Nov); Phase-1 code \hash{0578bce} and
notebook \hash{a2c5da2}, P10 marks \hash{abb1770} and result \hash{917ab3b}, Phase-1 start \hash{fec827e},
paper P10 and Phase-1 methods \hash{6e4175d} (21 Nov). Locked results: input lock of 13 November, hashes on Zenodo
(\href{https://doi.org/10.5281/zenodo.24206472}{doi:10.5281/zenodo.24206472}). Every number has a source
comment in this file. Claude (Anthropic) helped prepare this report.\par}
% src: GIT (all commits on origin/t3/fm-code-prep); PR l.10-11 (Zenodo)

\begin{thebibliography}{9}
\bibitem{yoo2021} T.\,K. Yoo, J.\,Y. Choi, H.\,K. Kim. Feasibility study to improve deep learning in OCT
diagnosis of rare retinal diseases with few-shot classification. \emph{Med.\ Biol.\ Eng.\ Comput.}
59:401--415, 2021.
\bibitem{sun2024} Sun, Wei, Wang, Zheng. \emph{Medical Image Analysis} 93:103102, 2024 (full title and DOI
not yet verified).
\end{thebibliography}

\end{document}
